\section{Marco teórico y trabajos relacionados}

Para fundamentar el desarrollo del Sistema de Carnetización Digital y Gestión de Eventos, se realizó una revisión sistemática de 20 artículos científicos publicados entre 2020 y 2025. Estos trabajos abordan problemas similares en la gestión de asistencia y verificación de identidad en contextos educativos e institucionales. El análisis se centró en tecnologías emergentes como códigos QR, RFID, blockchain y aprendizaje automático, evaluando su viabilidad, eficiencia y aceptación por parte de los usuarios. A continuación, se sintetizan los hallazgos clave, agrupados por temáticas principales.

\subsection{Sistemas de asistencia basados en códigos QR}

Varios estudios destacan el uso de códigos QR para simplificar el registro de asistencia, reemplazando métodos manuales ineficientes. Rajendran y Hamzah (2021) proponen un sistema de gestión de eventos con QR generado dinámicamente, integrando autenticación Firebase y escaneo móvil, logrando una reducción significativa en errores administrativos. En contextos educativos, Hamzah et al. (2021) implementan QR en la Universidad UiTM, automatizando cálculos de asistencia y envío de alertas, con calificaciones superiores a 4.0/5.0 en pruebas con 30 usuarios.

Deshmukh et al. (s.f.) desarrollan una aplicación web y móvil para instituciones educativas en India, utilizando PHP, MySQL y Android, permitiendo reportes en Excel. Sin embargo, identifican limitaciones en la sincronización en tiempo real. Hendrawan et al. (2025) realizan una revisión sistemática de 9 estudios, confirmando que los QR mejoran la precisión de datos y la gestión sostenible, aunque vulnerables al fraude si no se combinan con medidas adicionales como geofencing o IMEI.

En escenarios post-pandemia, Al Amin et al. (2024) integran QR con encriptación AES-128 para rastreo de contactos, priorizando seguridad y privacidad. Perin (2025) reutiliza una app de rastreo para asistencia en Bohol Island State University, obteniendo alta compatibilidad con dispositivos (4.08/5.0), pero con debilidades en estabilidad y seguridad.

\subsection{Integración de RFID para asistencia rápida}

El RFID emerge como tecnología complementaria para registros de alta velocidad. Mares et al. (2020) implementan un sistema en el Instituto Tecnológico Superior de Irapuato, utilizando Android Studio, Java y un lector RFID USB, logrando una reducción del 50% en tiempos de registro y 80% de efectividad en almacenamiento de datos. Sin embargo, sugieren validaciones adicionales para carnets autorizados.

\subsection{Sistemas híbridos y avanzados con múltiples capas de seguridad}

Para mitigar fraudes, varios trabajos proponen combinaciones tecnológicas. Siew et al. (2023) en i-CATS University College integran reconocimiento facial y QR dinámicos (cambiando cada 15 segundos), reduciendo tiempos de registro de 12 a 2 segundos y obteniendo 85% de preferencia por facial. Nwabuwe et al. (2023) usan QR dinámico, geofencing e IMEI, bloqueando exitosamente intentos de fraude simulados.

Dongre y Verma (2021) construyen un framework blockchain privado en Python para asistencia con QR, garantizando inmutabilidad de registros con tiempos de bloque inferiores a 0.017 segundos. Brunete y Cedazo (2018) proponen un sistema con QR y algoritmo CAOQTP para autenticación, combinando OTP por SMS con reordenamiento de imágenes CAPTCHA.

Zihan e Islam (2024) aplican aprendizaje automático (Mobile-ViT) para extracción y verificación de QR en aplicaciones musicales, logrando 99.28% de precisión en clasificación genuina/falsa. Matz et al. (2023) usan datos de apps móviles para predecir deserción estudiantil, combinando métricas académicas con comportamiento social, alcanzando 78% de precisión.

\subsection{Revisión sistemática y comparaciones metodológicas}

Razzaq y Ghayyur (2023) mapean 73 estudios sobre migración de monolíticos a microservicios, destacando beneficios en escalabilidad y agilidad, pero desafíos en comunicación distribuida. Kumar (s.f.) explora IA generativa en arquitectura de software, prediciendo reducciones del 60% en integración, aunque con riesgos de precisión.

Fitzpatrick y Storey (s.f.) analizan riesgos en integración continua, enfatizando la necesidad de supervisión humana. Beck (2003), Fowler y Foemmel (2006) y Martin (2008) fundamentan prácticas ágiles como TDD y refactorización para calidad.

Forsgren et al. (2020) y Dora (2021) correlacionan métricas DORA con rendimiento organizacional. Chen (2020) discute DevOps para entrega continua.

\subsection{Lecciones aprendidas y justificación de la solución propuesta}

De la revisión, emerge que los códigos QR estáticos ofrecen simplicidad y bajo costo, ideales para instituciones con recursos limitados, pero requieren capas adicionales para seguridad. Tecnologías avanzadas como blockchain o IA mejoran robustez, pero aumentan complejidad. El sistema propuesto adopta QR estáticos por universalidad y facilidad, con arquitectura preparada para expansiones futuras como RFID o geofencing, priorizando viabilidad sobre complejidad extrema.